%------------------------------------------------------------------------------%
\begin{question}
  Ortonormalize, pelo processo de Gram--Schmidt,
  \[
    \vec{v}_1 = (1,1,0)
    \qquad
    \vec{v}_2 = (1,0,1)
  \]
  e obtenha uma base ortonormal de
  \(
    \ger\{\vec{v}_1,\vec{v}_2\}
  \)
\end{question}

%------------------------------------------------------------------------------%
\begin{resolution}{Solução}
  \[
    \vec{u}_1
    \pause
    = \vec{v}_1
    \pause
    = (1, 1, 0)
  \]

  \pause
  \[
    \norm{\vec{u}_1}
    \pause
    = \sqrt{2}
  \]

  \pause
  \[
    \vec{e}_1
    \pause
    = \frac{1}{\sqrt{2}}(1,1,0)
  \]
\end{resolution}

%------------------------------------------------------------------------------%
\begin{resolution}{Solução}
\begin{columns}[t]
\column{0.3\textwidth}
  \begin{align*}
  \onslide<+->{\vec{u}_2}
  \onslide<+->{& = \vec{v}_2 - \frac{\vec{v}_2\cdot\vec{u}_1}{\norm{\vec{u}_1}^{2}} \; \vec{u}_1 \\}
  \onslide<+->{& = (1,0,1) - \frac{1}{2}(1,1,0) \\}
  \onslide<+->{& = \left( \frac{1}{2}, -\frac{1}{2}, 1 \right)}
  \end{align*}
\column{0.3\textwidth}
  \begin{align*}
    \onslide<+->{\norm{\vec{u}_2}^{2}}
    \onslide<+->{& = \frac{1}{4}+\frac{1}{4}+1\\}
    \onslide<+->{& = \frac{3}{2} \\}
    \onslide<+->{\norm{\vec{u}_2}}
    \onslide<+->{& = \frac{\sqrt{6}}{2}}
  \end{align*}
\column{0.3\textwidth}
  \begin{align*}
    \onslide<+->{\vec{e}_2}
    \onslide<+->{& = \frac{\vec{u}_2}{\norm{\vec{u}_2}}\\}
    \onslide<+->{& = \frac{1}{\sqrt{6}}(1,-1,2)}
  \end{align*}
\end{columns}
\end{resolution}

%------------------------------------------------------------------------------%
\begin{resolution}{Solução}
  \[
    \mathcal{B} =
    \left\{
      \frac{1}{\sqrt{2}}(1, 1, 0),
      \frac{1}{\sqrt{6}}(1, -1, 2)
    \right\}
  \]
\end{resolution}

%------------------------------------------------------------------------------%
\endinput
